\chapter{Literature Review and Related Work}
\label{chap:relatedworks}

\section{Competitor Analysis}
\label{section:competitor-analysis}

A competitor analysis was conducted to identify the strengths and
limitations of existing systems in the skating performance analysis
domain. Table~\ref{tab:competitor} summarizes each system across
three dimensions: focus area and key limitations.

\begin{table}[h]
    \centering
    \renewcommand{\arraystretch}{1.3}
    \caption{Competitor Analysis of Skating Jump Analysis Systems}
    \label{tab:competitor}
    \begin{tabular}{|l|l|l|}
        \hline
        \textbf{System}
            & \textbf{Focus}
            & \textbf{Limitations} \\
        \hline
        Athlix
            & Post-match stats
            & Offline only (5 to 10s latency) \\
        \hline
        Dartfish
            & Training video markup
            & Manual labeling, no AI \\
        \hline
        YourSkatingCoach
            & Air-time measurement
            & Single metric, 65\% accuracy \\
        \hline
        FigureFlow
            & Live jump recognition
            & N/A \\
        \hline
    \end{tabular}
\end{table}

\textbf{Dartfish} is a professional video analysis platform used by Olympic coaches for technical breakdown of athlete performances. It enables manual tagging of key moments, slow-motion review, and detailed annotations, making it valuable for post-competition analysis. However, Dartfish requires extensive human expertise for tagging each jump and cannot provide real-time predictions during live broadcasts, failing to address FigureFlow's goal of automated, low-latency jump identification for casual viewers.

\textbf{Athlix} delivers AI-powered injury prevention and performance optimization through posture analysis and movement tracking. While innovative for athlete health monitoring across webcam-based assessments, it focuses on general sports recovery rather than specialized figure skating jump classification. Athlix lacks the frame-by-frame takeoff recognition needed for live jump identification in competitions.

The analysis reveals a clear gap in the existing landscape: no current
system combines real-time processing, automated AI classification, and
multi-feature pose analysis into a single deployment-ready pipeline.
FigureFlow is designed to address all three of these limitations
simultaneously.


\section{Literature Review}
\label{section:literature-review}

This section reviews academic papers and technical reports relevant
to FigureFlow. For each source, the key findings are summarized and
connected to specific design and implementation decisions in this project.

\textbf{Pose-based action recognition.}
Ngu et al.~\cite{ieee2023_pose} evaluated MediaPipe Pose combined with
an LSTM classifier on the UTD-MHAD benchmark dataset, achieving 92\%
classification accuracy. Their results demonstrate that skeletal keypoint
sequences generalize well across subjects and capture motion structure
that RGB pixel data cannot reliably encode. This finding directly informs
the decision to use MediaPipe as the primary feature extractor in
FigureFlow, replacing raw video frames with normalized pose sequences as
input to the classifier.

\textbf{Temporal sequence modeling.}
Chen et al.~\cite{cvpr2024_skating} presented a comparative study of
sequence models on skating-specific pose data at the CVPR 2024 workshop,
showing that a Bidirectional LSTM (Bi-LSTM) outperforms CNN-based
classifiers by 8\% on jump phase segmentation tasks. The bidirectional
formulation captures both forward and backward temporal context, which is
critical for distinguishing jumps whose takeoff and landing characteristics
overlap in time. This justifies the adoption of an LSTM-based temporal
model in our classification pipeline.

\textbf{Real-time object detection.}
Jocher et al.~\cite{ultralytics_yolov8} introduced YOLOv8, which
establishes the current state-of-the-art for real-time object detection
across the COCO and MOT benchmarks. In FigureFlow, YOLOv8 is employed
in the first stage of the pipeline to localize and track the skater
within each frame before pose estimation is applied. This reduces noise
caused by background subjects and ensures that pose keypoints correspond
exclusively to the athlete of interest.

Taken together, these works validate the three-stage architecture of
FigureFlow: YOLOv8 for detection, MediaPipe for pose extraction, and
LSTM for temporal classification. Each design choice is grounded in
empirical evidence from the reviewed literature, and the combination
addresses the limitations identified in the competitor analysis,
particularly the inability of existing tools to perform fully automated,
real-time jump recognition.